\documentclass[11pt]{article}

\usepackage{amsmath,amssymb,amsthm}
\usepackage{geometry}
\usepackage{bm}
\usepackage[framemethod=TikZ]{mdframed}

% Table of Contents depth
\setcounter{tocdepth}{2}      % sections and subsections
\setcounter{secnumdepth}{3}   % number down to subsubsections

% Hyperlinks (load last)
\usepackage[
colorlinks=true,
linkcolor=blue,
citecolor=blue,
urlcolor=blue
]{hyperref}
\usepackage{bookmark}

\geometry{margin=1in}

\title{Expected Value}
\author{}
\date{}

\begin{document}
	
	\maketitle
	\tableofcontents
	
	
	\section{Definition of Expectation (Discrete and Continuous)}
	\begin{mdframed}[
		linewidth=0.8pt,
		roundcorner=6pt,
		backgroundcolor=gray!5,
		linecolor=gray!60,
		innertopmargin=10pt,
		innerbottommargin=10pt,
		innerleftmargin=10pt,
		innerrightmargin=10pt
		]
		
		\subsection*{Definition of the Expected Value}
		
		The expected value (or \emph{expectation}) of a random variable is one of the
		most fundamental concepts in probability theory and statistics. It
		represents the average or center of the probability distribution and can be
		interpreted as the value that would be obtained by averaging the outcomes of
		the random variable over an indefinitely large number of independent
		repetitions of the experiment.
		
		More formally, the expected value is a weighted average of all possible
		values of the random variable, where each value is weighted according to its
		probability of occurring.
		
		The precise definition depends on whether the random variable is discrete or
		continuous.
		
		\bigskip
		
		\paragraph{Discrete Random Variables}
		
		Suppose that $X$ is a discrete random variable taking values
		
		\[
		x_1,x_2,x_3,\ldots
		\]
		
		with probability mass function
		
		\[
		P(X=x_i)=p_i,
		\qquad
		i=1,2,\ldots,
		\]
		
		where
		
		\[
		p_i\ge0,
		\qquad
		\sum_{i=1}^{\infty}p_i=1.
		\]
		
		If the series
		
		\[
		\sum_{i=1}^{\infty}|x_i|p_i
		\]
		
		is finite, then the expected value of $X$ is defined as
		
		\[
		\boxed{
			E[X]
			=
			\sum_{i=1}^{\infty}x_ip_i.
		}
		\]
		
		Equivalently,
		
		\[
		\boxed{
			E[X]
			=
			\sum_x x\,P(X=x),
		}
		\]
		
		where the summation extends over all possible values of $X$.
		
		Observe that each possible value contributes proportionally to its
		probability. Values with larger probabilities contribute more heavily to the
		expected value than values with smaller probabilities.
		
		\bigskip
		
		\paragraph{Example}
		
		Consider the random variable
		
		\[
		X=
		\begin{cases}
			1,&P=\dfrac14,\\[0.2cm]
			2,&P=\dfrac12,\\[0.2cm]
			5,&P=\dfrac14.
		\end{cases}
		\]
		
		Then
		
		\[
		\begin{aligned}
			E[X]
			&=
			1\left(\frac14\right)
			+
			2\left(\frac12\right)
			+
			5\left(\frac14\right)
			\\[0.3cm]
			&=
			\frac14
			+
			1
			+
			\frac54
			\\[0.3cm]
			&=
			\frac{10}{4}
			\\[0.2cm]
			&=
			2.5.
		\end{aligned}
		\]
		
		Notice that the expected value need not be one of the possible values of the
		random variable. Although $X$ can only take the values
		
		\[
		1,\;2,\;\text{or}\;5,
		\]
		
		its expected value equals
		
		\[
		2.5.
		\]
		
		The expectation should therefore be interpreted as a long-run average rather
		than as a value that the random variable must actually attain.
		
		\bigskip
		
		\paragraph{Continuous Random Variables}
		
		Suppose now that $X$ is a continuous random variable having probability
		density function
		
		\[
		f_X(x),
		\]
		
		where
		
		\[
		f_X(x)\ge0,
		\]
		
		and
		
		\[
		\int_{-\infty}^{\infty}f_X(x)\,dx=1.
		\]
		
		If
		
		\[
		\int_{-\infty}^{\infty}|x|f_X(x)\,dx<\infty,
		\]
		
		then the expected value of $X$ is defined as
		
		\[
		\boxed{
			E[X]
			=
			\int_{-\infty}^{\infty}
			x\,f_X(x)\,dx.
		}
		\]
		
		This definition is the continuous analogue of the discrete case.
		
		Instead of summing over all possible values, we integrate over the entire
		real line.
		
		Again, each possible value of $X$ is weighted by the probability density at
		that point.
		
		\bigskip
		
		\paragraph{Example}
		
		Suppose
		
		\[
		X\sim U(0,1),
		\]
		
		so that
		
		\[
		f_X(x)=
		\begin{cases}
			1,&0\le x\le1,\\
			0,&\text{otherwise}.
		\end{cases}
		\]
		
		The expected value is
		
		\[
		\begin{aligned}
			E[X]
			&=
			\int_{-\infty}^{\infty}
			x\,f_X(x)\,dx
			\\[0.3cm]
			&=
			\int_0^1x\,dx
			\\[0.3cm]
			&=
			\left[
			\frac{x^2}{2}
			\right]_0^1
			\\[0.3cm]
			&=
			\frac12.
		\end{aligned}
		\]
		
		Thus,
		
		\[
		\boxed{
			E[X]=\frac12.
		}
		\]
		
		As expected, the center of the uniform distribution over the interval
		$[0,1]$ is precisely its midpoint.
		
		\bigskip
		
		\paragraph{Existence of the Expected Value}
		
		The expected value is not always finite.
		
		For example, if
		
		\[
		P(X=n)
		=
		\frac{6}{\pi^2n^2},
		\qquad
		n=1,2,\ldots,
		\]
		
		then
		
		\[
		\sum_{n=1}^{\infty}P(X=n)=1,
		\]
		
		so this is a valid probability distribution.
		
		However,
		
		\[
		E[X]
		=
		\sum_{n=1}^{\infty}
		n
		\cdot
		\frac{6}{\pi^2n^2}
		=
		\frac6{\pi^2}
		\sum_{n=1}^{\infty}
		\frac1n.
		\]
		
		Since the harmonic series diverges,
		
		\[
		\sum_{n=1}^{\infty}\frac1n=\infty,
		\]
		
		it follows that
		
		\[
		\boxed{
			E[X]=\infty.
		}
		\]
		
		For this reason, the definition of expectation always requires the
		additional condition
		
		\[
		\sum|x|P(X=x)<\infty
		\]
		
		or
		
		\[
		\int|x|f(x)\,dx<\infty,
		\]
		
		which guarantees that the expectation is well defined.
		
		\bigskip
		
		\paragraph{Unified Definition}
		
		The discrete and continuous definitions can be viewed as special cases of a
		single principle.
		
		The expected value is obtained by averaging a quantity with respect to the
		probability distribution of the random variable.
		
		Symbolically,
		
		\[
		\boxed{
			E[X]
			=
			\int x\,dP,
		}
		\]
		
		where the integral is understood in the sense of probability theory.
		
		This unified viewpoint allows all subsequent properties of expectation to be
		proved simultaneously for both discrete and continuous random variables.
		
	\end{mdframed}
	
	
	\section{Expectation of a Function (LOTUS)}
	\begin{mdframed}[
		linewidth=0.8pt,
		roundcorner=6pt,
		backgroundcolor=gray!5,
		linecolor=gray!60,
		innertopmargin=10pt,
		innerbottommargin=10pt,
		innerleftmargin=10pt,
		innerrightmargin=10pt
		]
		
		\subsection*{Expectation of a Function (The Law of the Unconscious Statistician)}
		
		In many applications, we are interested not only in the expected value of a
		random variable $X$, but also in the expected value of some function of that
		random variable,
		
		\[
		g(X),
		\]
		
		where $g$ is a real-valued function.
		
		For example, we may wish to compute
		
		\[
		E[X^2],
		\qquad
		E[e^X],
		\qquad
		E[\ln X],
		\qquad
		E[\sin X],
		\]
		
		or any other transformation of the original random variable.
		
		A natural question therefore arises:
		
		\begin{center}
			\emph{Must we first determine the probability distribution of $g(X)$ before
				computing its expected value?}
		\end{center}
		
		Fortunately, the answer is no.
		
		The expected value of a function of a random variable can be computed
		directly from the probability distribution of the original random variable.
		This remarkable result is known as the \textbf{Law of the Unconscious
			Statistician (LOTUS).}
		
		\bigskip
		
		\paragraph{Discrete Random Variables}
		
		Suppose that $X$ is a discrete random variable taking values
		
		\[
		x_1,x_2,\ldots
		\]
		
		with probability mass function
		
		\[
		P(X=x_i)=p_i.
		\]
		
		The random variable
		
		\[
		g(X)
		\]
		
		takes the values
		
		\[
		g(x_1),g(x_2),\ldots,
		\]
		
		possibly with repeated values.
		
		Instead of determining the probability distribution of $g(X)$, we compute
		its expectation directly.
		
		The expected value is defined by
		
		\[
		\boxed{
			E[g(X)]
			=
			\sum_x
			g(x)P(X=x),
		}
		\]
		
		provided that
		
		\[
		\sum_x
		|g(x)|P(X=x)
		<
		\infty.
		\]
		
		Observe that this formula is identical to the definition of expectation,
		except that the values $x$ have simply been replaced by $g(x)$.
		
		\bigskip
		
		\paragraph{Example}
		
		Suppose
		
		\[
		P(X=1)=0.2,
		\qquad
		P(X=2)=0.5,
		\qquad
		P(X=4)=0.3.
		\]
		
		Compute
		
		\[
		E[X^2].
		\]
		
		Rather than finding the distribution of $X^2$, we apply LOTUS directly,
		
		\[
		\begin{aligned}
			E[X^2]
			&=
			1^2(0.2)
			+
			2^2(0.5)
			+
			4^2(0.3)
			\\[0.3cm]
			&=
			0.2
			+
			2
			+
			4.8
			\\[0.3cm]
			&=
			7.
		\end{aligned}
		\]
		
		Notice that the probability distribution of $X^2$ never had to be computed.
		
		\bigskip
		
		\paragraph{Continuous Random Variables}
		
		Suppose now that $X$ is a continuous random variable with probability
		density function
		
		\[
		f_X(x).
		\]
		
		Then the expected value of $g(X)$ is
		
		\[
		\boxed{
			E[g(X)]
			=
			\int_{-\infty}^{\infty}
			g(x)f_X(x)\,dx,
		}
		\]
		
		provided that
		
		\[
		\int_{-\infty}^{\infty}
		|g(x)|f_X(x)\,dx
		<
		\infty.
		\]
		
		Again, the distribution of $g(X)$ is not required.
		
		Instead, the function $g$ is evaluated at every possible value of $X$, and
		the result is averaged using the density of $X$.
		
		\bigskip
		
		\paragraph{Example}
		
		Suppose
		
		\[
		X\sim U(0,1),
		\]
		
		whose density is
		
		\[
		f(x)=
		\begin{cases}
			1,&0\le x\le1,\\
			0,&\text{otherwise}.
		\end{cases}
		\]
		
		Compute
		
		\[
		E[X^2].
		\]
		
		Using LOTUS,
		
		\[
		\begin{aligned}
			E[X^2]
			&=
			\int_0^1x^2\,dx
			\\[0.3cm]
			&=
			\left[
			\frac{x^3}{3}
			\right]_0^1
			\\[0.3cm]
			&=
			\frac13.
		\end{aligned}
		\]
		
		Thus,
		
		\[
		\boxed{
			E[X^2]=\frac13.
		}
		\]
		
		Again, there was no need to determine the probability density of the random
		variable $X^2$.
		
		\bigskip
		
		\paragraph{Why Does LOTUS Work? (Discrete Case)}
		
		Suppose that
		
		\[
		Y=g(X).
		\]
		
		By definition,
		
		\[
		E[Y]
		=
		\sum_y
		yP(Y=y).
		\]
		
		Since
		
		\[
		Y=g(X),
		\]
		
		the event
		
		\[
		Y=y
		\]
		
		occurs whenever
		
		\[
		g(X)=y.
		\]
		
		Hence,
		
		\[
		P(Y=y)
		=
		P(g(X)=y).
		\]
		
		Substituting,
		
		\[
		E[Y]
		=
		\sum_y
		yP(g(X)=y).
		\]
		
		Now collect together all values of $X$ that satisfy
		
		\[
		g(x)=y.
		\]
		
		Since
		
		\[
		P(g(X)=y)
		=
		\sum_{x:g(x)=y}
		P(X=x),
		\]
		
		we obtain
		
		\[
		\begin{aligned}
			E[Y]
			&=
			\sum_y
			y
			\sum_{x:g(x)=y}
			P(X=x)
			\\[0.3cm]
			&=
			\sum_x
			g(x)P(X=x).
		\end{aligned}
		\]
		
		Therefore,
		
		\[
		\boxed{
			E[g(X)]
			=
			\sum_x
			g(x)P(X=x).
		}
		\]
		
		Thus, the probability distribution of $g(X)$ never needs to be computed
		explicitly.
		
		\bigskip
		
		\paragraph{Why Is It Called the ``Law of the Unconscious Statistician''?}
		
		The name originates from the fact that many people compute
		
		\[
		E[g(X)]
		=
		\sum_xg(x)P(X=x)
		\]
		
		or
		
		\[
		E[g(X)]
		=
		\int g(x)f(x)\,dx
		\]
		
		without consciously realizing that they are using a theorem.
		
		Instead of first determining the probability distribution of $g(X)$, they
		simply replace the variable $x$ in the expectation formula by $g(x)$ and
		proceed with the calculation.
		
		LOTUS guarantees that this procedure is mathematically correct.
		
		\bigskip
		
		\paragraph{Special Cases}
		
		Several important expectations are immediate consequences of LOTUS.
		
		Choosing
		
		\[
		g(x)=x,
		\]
		
		gives
		
		\[
		E[g(X)]
		=
		E[X].
		\]
		
		Choosing
		
		\[
		g(x)=x^2,
		\]
		
		gives
		
		\[
		E[X^2].
		\]
		
		Choosing
		
		\[
		g(x)=|x|,
		\]
		
		gives
		
		\[
		E[|X|].
		\]
		
		Choosing
		
		\[
		g(x)=e^x,
		\]
		
		gives
		
		\[
		E[e^X].
		\]
		
		Thus, virtually every expectation encountered in probability and statistics
		is an application of LOTUS.
		
		\bigskip
		
		\paragraph{Conclusion}
		
		The Law of the Unconscious Statistician provides a powerful computational
		tool.
		
		Rather than determining the probability distribution of a transformed
		random variable $g(X)$, one computes its expectation directly from the
		probability distribution of the original random variable.
		
		For a discrete random variable,
		
		\[
		\boxed{
			E[g(X)]
			=
			\sum_xg(x)P(X=x),
		}
		\]
		
		while for a continuous random variable,
		
		\[
		\boxed{
			E[g(X)]
			=
			\int_{-\infty}^{\infty}
			g(x)f_X(x)\,dx.
		}
		\]
		
		Nearly every property of expectation developed in the following sections is
		a direct consequence of these two formulas.
		
	\end{mdframed}
	
	\section{Expectation of a Constant}
	\begin{mdframed}[
		linewidth=0.8pt,
		roundcorner=6pt,
		backgroundcolor=gray!5,
		linecolor=gray!60,
		innertopmargin=10pt,
		innerbottommargin=10pt,
		innerleftmargin=10pt,
		innerrightmargin=10pt
		]
		
		\subsection*{Expectation of a Constant}
		
		One of the simplest yet most fundamental properties of the expected value is
		that the expectation of a constant random variable is the constant itself.
		
		Intuitively, if a random variable always takes the same value, then there is
		no randomness to average. Consequently, its expected value must equal that
		constant.
		
		We now prove this result separately for discrete and continuous random
		variables.
		
		\bigskip
		
		\paragraph{Theorem}
		
		Let
		
		\[
		X=c,
		\]
		
		where \(c\) is a fixed real number.
		
		Then
		
		\[
		\boxed{
			E[c]=c.
		}
		\]
		
		\bigskip
		
		\paragraph{Proof for a Discrete Random Variable}
		
		Suppose that the random variable always takes the value
		
		\[
		c.
		\]
		
		Its probability mass function is therefore
		
		\[
		P(X=c)=1,
		\]
		
		and
		
		\[
		P(X=x)=0,
		\qquad
		x\neq c.
		\]
		
		By the definition of expectation,
		
		\[
		E[X]
		=
		\sum_x
		xP(X=x).
		\]
		
		Since every probability is zero except at \(x=c\),
		
		\[
		\begin{aligned}
			E[X]
			&=
			cP(X=c)
			+
			\sum_{x\neq c}
			xP(X=x)
			\\[0.3cm]
			&=
			c(1)
			+
			0
			\\[0.3cm]
			&=
			c.
		\end{aligned}
		\]
		
		Hence,
		
		\[
		\boxed{
			E[c]=c.
		}
		\]
		
		\bigskip
		
		\paragraph{Proof for a Continuous Random Variable}
		
		Although a constant random variable is technically not continuous (its
		distribution is concentrated at a single point), it is instructive to see
		how the same conclusion follows formally.
		
		Suppose that
		
		\[
		X=c.
		\]
		
		Its probability distribution places all its probability mass at the point
		\(c\). In measure-theoretic language,
		
		\[
		f_X(x)
		=
		\delta(x-c),
		\]
		
		where \(\delta\) denotes the Dirac delta distribution.
		
		Applying the definition of expectation,
		
		\[
		\begin{aligned}
			E[X]
			&=
			\int_{-\infty}^{\infty}
			x\,\delta(x-c)\,dx
			\\[0.3cm]
			&=
			c.
		\end{aligned}
		\]
		
		Thus,
		
		\[
		\boxed{
			E[c]=c.
		}
		\]
		
		\medskip
		
		For an introductory course, the discrete proof is sufficient. The
		measure-theoretic formulation above simply confirms that the same property
		holds in complete generality.
		
		\bigskip
		
		\paragraph{Alternative Proof Using LOTUS}
		
		The result also follows immediately from the Law of the Unconscious
		Statistician.
		
		Let
		
		\[
		g(x)=c,
		\]
		
		that is, \(g\) is the constant function.
		
		Applying LOTUS,
		
		\[
		E[g(X)]
		=
		\sum_x
		g(x)P(X=x)
		\]
		
		for a discrete random variable, or
		
		\[
		E[g(X)]
		=
		\int_{-\infty}^{\infty}
		g(x)f_X(x)\,dx
		\]
		
		for a continuous random variable.
		
		Since
		
		\[
		g(x)=c
		\]
		
		for every value of \(x\),
		
		\[
		\begin{aligned}
			E[g(X)]
			&=
			\sum_x
			cP(X=x)
			\\[0.3cm]
			&=
			c
			\sum_x
			P(X=x)
			\\[0.3cm]
			&=
			c(1)
			\\[0.3cm]
			&=
			c.
		\end{aligned}
		\]
		
		Similarly, in the continuous case,
		
		\[
		\begin{aligned}
			E[g(X)]
			&=
			\int_{-\infty}^{\infty}
			cf_X(x)\,dx
			\\[0.3cm]
			&=
			c
			\int_{-\infty}^{\infty}
			f_X(x)\,dx
			\\[0.3cm]
			&=
			c(1)
			\\[0.3cm]
			&=
			c.
		\end{aligned}
		\]
		
		Therefore,
		
		\[
		\boxed{
			E[c]=c.
		}
		\]
		
		\bigskip
		
		\paragraph{Example}
		
		Suppose a machine always produces an object weighing exactly
		
		\[
		5\text{ kg}.
		\]
		
		The corresponding random variable satisfies
		
		\[
		X=5.
		\]
		
		Its expectation is therefore
		
		\[
		E[X]=5.
		\]
		
		Although expectation is often interpreted as a long-run average, in this
		case every observation equals \(5\), so the average is trivially equal to
		\(5\).
		
		\bigskip
		
		\paragraph{Importance of the Result}
		
		Although elementary, this property is fundamental because it serves as one
		of the building blocks of the linearity of expectation.
		
		Indeed, when proving
		
		\[
		E[aX+b]
		=
		aE[X]+b,
		\]
		
		the term
		
		\[
		E[b]
		\]
		
		appears naturally. The present theorem immediately yields
		
		\[
		E[b]=b,
		\]
		
		allowing the linearity formula to be established.
		
		Thus, the expectation of a constant is not merely a simple observation; it
		is an essential ingredient in many subsequent results concerning expected
		values.
		
	\end{mdframed}
	
	\section{Linearity of Expectation}
	\begin{mdframed}[
		linewidth=0.8pt,
		roundcorner=6pt,
		backgroundcolor=gray!5,
		linecolor=gray!60,
		innertopmargin=10pt,
		innerbottommargin=10pt,
		innerleftmargin=10pt,
		innerrightmargin=10pt
		]
		
		\subsection*{Linearity of Expectation}
		
		The most important property of the expected value is its \emph{linearity}.
		Unlike many other statistical quantities, such as the variance,
		
		\[
		\operatorname{Var}(X+Y)
		\neq
		\operatorname{Var}(X)+\operatorname{Var}(Y)
		\]
		
		in general.
		
		Expectation behaves much more simply.
		
		The expectation of a linear combination of random variables is the
		corresponding linear combination of their expectations.
		
		Remarkably, this property holds whether or not the random variables are
		independent.
		
		This fact makes expectation one of the most useful operators in probability
		and statistics.
		
		\bigskip
		
		\paragraph{Theorem}
		
		Let
		
		\[
		X
		\quad\text{and}\quad
		Y
		\]
		
		be random variables whose expectations exist.
		
		Let
		
		\[
		a,b\in\mathbb R
		\]
		
		be constants.
		
		Then
		
		\[
		\boxed{
			E[aX+bY]
			=
			aE[X]+bE[Y].
		}
		\]
		
		More generally, if
		
		\[
		X_1,\ldots,X_n
		\]
		
		are random variables with finite expectations, then
		
		\[
		\boxed{
			E\!\left[
			\sum_{i=1}^{n}a_iX_i
			\right]
			=
			\sum_{i=1}^{n}
			a_iE[X_i].
		}
		\]
		
		Notice that no assumption of independence is required.
		
		\bigskip
		
		\paragraph{Proof for Discrete Random Variables}
		
		Suppose that
		
		\[
		X \qquad \text{and} \qquad Y
		\]
	
		
		have joint probability mass function
		
		\[
		P(X=x,Y=y).
		\]
		
		By the definition of expectation,
		
		\[
		E[aX+bY]
		=
		\sum_x
		\sum_y
		(a x+b y)
		P(X=x,Y=y).
		\]
		
		Using the distributive property,
		
		\[
		\begin{aligned}
			E[aX+bY]
			&=
			\sum_x
			\sum_y
			\left(
			ax+by
			\right)
			P(X=x,Y=y)
			\\[0.3cm]
			&=
			\sum_x
			\sum_y
			axP(X=x,Y=y)
			\\
			&\qquad
			+
			\sum_x
			\sum_y
			byP(X=x,Y=y).
		\end{aligned}
		\]
		
		Since \(a\) and \(b\) are constants,
		
		\[
		\begin{aligned}
			E[aX+bY]
			&=
			a
			\sum_x
			\sum_y
			xP(X=x,Y=y)
			\\
			&\qquad
			+
			b
			\sum_x
			\sum_y
			yP(X=x,Y=y).
		\end{aligned}
		\]
		
		Now consider the first double summation.
		
		Using the law of total probability,
		
		\[
		\sum_y
		P(X=x,Y=y)
		=
		P(X=x).
		\]
		
		Therefore,
		
		\[
		\begin{aligned}
			\sum_x
			\sum_y
			xP(X=x,Y=y)
			&=
			\sum_x
			x
			\left(
			\sum_y
			P(X=x,Y=y)
			\right)
			\\[0.3cm]
			&=
			\sum_x
			xP(X=x)
			\\[0.3cm]
			&=
			E[X].
		\end{aligned}
		\]
		
		Similarly,
		
		\[
		\begin{aligned}
			\sum_x
			\sum_y
			yP(X=x,Y=y)
			&=
			\sum_y
			y
			\left(
			\sum_x
			P(X=x,Y=y)
			\right)
			\\[0.3cm]
			&=
			\sum_y
			yP(Y=y)
			\\[0.3cm]
			&=
			E[Y].
		\end{aligned}
		\]
		
		Substituting these results into the previous expression,
		
		\[
		\begin{aligned}
			E[aX+bY]
			&=
			aE[X]
			+
			bE[Y].
		\end{aligned}
		\]
		
		Hence,
		
		\[
		\boxed{
			E[aX+bY]
			=
			aE[X]+bE[Y].
		}
		\]
		
		\bigskip
		
		\paragraph{Proof for Continuous Random Variables}
		
		Suppose that
		
		\[
		X \qquad \text{and} \qquad Y
		\]
		
		have joint probability density function
		
		\[
		f(x,y).
		\]
		
		The expectation of
		
		\[
		aX+bY
		\]
		
		is
		
		\[
		E[aX+bY]
		=
		\int_{-\infty}^{\infty}
		\int_{-\infty}^{\infty}
		(ax+by)
		f(x,y)
		\,dx\,dy.
		\]
		
		Using the linearity of integration,
		
		\[
		\begin{aligned}
			E[aX+bY]
			&=
			a
			\int\!\!\!\int
			x
			f(x,y)
			\,dx\,dy
			\\
			&\qquad
			+
			b
			\int\!\!\!\int
			y
			f(x,y)
			\,dx\,dy.
		\end{aligned}
		\]
		
		\paragraph{Evaluating the First Integral}
		
		Let us carefully study the first double integral,
		
		\[
		\int_{-\infty}^{\infty}
		\int_{-\infty}^{\infty}
		x\,f(x,y)\,dx\,dy.
		\]
		
		Notice that the integrand is
		
		\[
		x\,f(x,y),
		\]
		
		where the variable \(x\) depends only on \(x\), while the joint density
		\(f(x,y)\) depends on both variables.
		
		Our objective is to show that this double integral is precisely the
		expectation of the random variable \(X\).
		
		\bigskip
		
		The first step is to change the order of integration.
		
		Since the joint density is nonnegative and integrable, Fubini's Theorem
		allows us to interchange the order of integration:
		
		\[
		\begin{aligned}
			\int_{-\infty}^{\infty}
			\int_{-\infty}^{\infty}
			x\,f(x,y)
			\,dx\,dy
			&=
			\int_{-\infty}^{\infty}
			\left(
			\int_{-\infty}^{\infty}
			x\,f(x,y)
			\,dy
			\right)
			dx.
		\end{aligned}
		\]
		
		Notice that the integration is now performed with respect to \(y\) first.
		
		\bigskip
		
		Next, observe that during the inner integration,
		
		\[
		\int_{-\infty}^{\infty}
		x\,f(x,y)\,dy,
		\]
		
		the variable \(x\) is treated as a constant because the integration variable
		is \(y\).
		
		Therefore,
		
		\[
		\begin{aligned}
			\int_{-\infty}^{\infty}
			x\,f(x,y)\,dy
			&=
			x
			\int_{-\infty}^{\infty}
			f(x,y)\,dy.
		\end{aligned}
		\]
		
		Substituting this into the previous expression gives
		
		\[
		\begin{aligned}
			\int_{-\infty}^{\infty}
			\int_{-\infty}^{\infty}
			x\,f(x,y)
			\,dx\,dy
			&=
			\int_{-\infty}^{\infty}
			x
			\left(
			\int_{-\infty}^{\infty}
			f(x,y)\,dy
			\right)
			dx.
		\end{aligned}
		\]
		
		\bigskip
		
		\paragraph{Deriving the Marginal Density}
		
		The quantity
		
		\[
		f(x,y)
		\]
		
		is the \emph{joint probability density function} of the random variables
		\(X\) and \(Y\).
		
		It describes how probability is distributed over the two-dimensional
		\((x,y)\)-plane.
		
		Suppose we are interested only in the random variable \(X\). In other words,
		we no longer care about the value of \(Y\).
		
		To obtain the probability density of \(X\) alone, we must account for
		\emph{every possible value} that \(Y\) can take.
		
		This is accomplished by integrating the joint density over the entire range
		of \(Y\):
		
		\[
		\boxed{
			f_X(x)
			=
			\int_{-\infty}^{\infty}
			f(x,y)\,dy.
		}
		\]
		
		This function is called the \emph{marginal density} of \(X\).
		
		\bigskip
		
		\paragraph{Why Does This Formula Work?}
		
		Recall that for continuous random variables,
		
		\[
		P\!\left(
		a\le X\le b,\,
		c\le Y\le d
		\right)
		=
		\int_a^b
		\int_c^d
		f(x,y)
		\,dy\,dx.
		\]
		
		Now suppose we wish to compute the probability that
		
		\[
		X
		\]
		
		lies between \(a\) and \(b\), regardless of the value of \(Y\).
		
		Since \(Y\) may take any value on the real line, we integrate over all
		possible values of \(Y\):
		
		\[
		\begin{aligned}
			P(a\le X\le b)
			&=
			\int_a^b
			\int_{-\infty}^{\infty}
			f(x,y)
			\,dy\,dx.
		\end{aligned}
		\]
		
		On the other hand, by the definition of the marginal density,
		
		\[
		P(a\le X\le b)
		=
		\int_a^b
		f_X(x)\,dx.
		\]
		
		Since these two expressions represent exactly the same probability for every
		interval \([a,b]\),
		
		\[
		\int_a^b
		f_X(x)\,dx
		=
		\int_a^b
		\left(
		\int_{-\infty}^{\infty}
		f(x,y)\,dy
		\right)
		dx.
		\]
		
		The Fundamental Theorem of Calculus then implies that the integrands must be
		equal almost everywhere. Therefore,
		
		\[
		\boxed{
			f_X(x)
			=
			\int_{-\infty}^{\infty}
			f(x,y)\,dy.
		}
		\]
		
		\bigskip
		
		\paragraph{Geometric Interpretation}
		
		The joint density
		
		\[
		f(x,y)
		\]
		
		may be visualized as a surface lying above the \((x,y)\)-plane.
		
		For a fixed value of \(x\), the function
		
		\[
		y\mapsto f(x,y)
		\]
		
		is a vertical slice through this surface.
		
		Integrating
		
		\[
		f(x,y)
		\]
		
		with respect to \(y\) computes the area under this slice.
		
		That area is precisely the value of the marginal density
		
		\[
		f_X(x).
		\]
		
		Thus,
		
		\[
		f_X(x)
		\]
		
		measures the total probability density associated with the value \(x\),
		after accounting for every possible value of the second variable \(Y\).
		
		\bigskip
		
		\paragraph{Example}
		
		Consider the joint density
		
		\[
		f(x,y)
		=
		\begin{cases}
			2,
			&
			0<x<y<1,
			\\[0.2cm]
			0,
			&
			\text{otherwise.}
		\end{cases}
		\]
		
		To obtain the marginal density of \(X\), we integrate over all admissible
		values of \(Y\).
		
		Since the condition
		
		\[
		0<x<y<1
		\]
		
		implies that
		
		\[
		y
		\in
		(x,1),
		\]
		
		we have
		
		\[
		\begin{aligned}
			f_X(x)
			&=
			\int_x^1
			2
			\,dy
			\\[0.3cm]
			&=
			2(1-x),
			\qquad
			0<x<1.
		\end{aligned}
		\]
		
		Therefore,
		
		\[
		\boxed{
			f_X(x)
			=
			\begin{cases}
				2(1-x),
				&
				0<x<1,
				\\[0.2cm]
				0,
				&
				\text{otherwise.}
			\end{cases}
		}
		\]
		
		Notice that the variable \(Y\) has completely disappeared. All of its
		possible values have been "summed out" through integration, leaving a
		density that depends only on \(X\).
		
		Intuitively, this integral "adds up" the probability density corresponding
		to every possible value of \(Y\), leaving only the probability density of
		\(X\).
		
		Replacing the inner integral by the marginal density,
		
		\[
		\begin{aligned}
			\int_{-\infty}^{\infty}
			\int_{-\infty}^{\infty}
			x\,f(x,y)
			\,dx\,dy
			&=
			\int_{-\infty}^{\infty}
			x
			f_X(x)
			\,dx.
		\end{aligned}
		\]
		
		\bigskip
		
		Finally, recall the definition of the expectation of a continuous random
		variable:
		
		\[
		E[X]
		=
		\int_{-\infty}^{\infty}
		x
		f_X(x)
		\,dx.
		\]
		
		Therefore,
		
		\[
		\boxed{
			\int_{-\infty}^{\infty}
			\int_{-\infty}^{\infty}
			x\,f(x,y)
			\,dx\,dy
			=
			E[X].
		}
		\]
		
		The same argument applied to the second integral yields
		
		\[
		\boxed{
			\int_{-\infty}^{\infty}
			\int_{-\infty}^{\infty}
			y\,f(x,y)
			\,dx\,dy
			=
			E[Y].
		}
		\]
		
		Consequently,
		
		\[
		E[aX+bY]
		=
		aE[X]
		+
		bE[Y].
		\]
		
		\bigskip
		
		\paragraph{Generalization}
		
		The previous result extends naturally to any finite number of random
		variables.
		
		Indeed,
		
		\[
		\begin{aligned}
			E[a_1X_1+\cdots+a_nX_n]
			&=
			E[a_1X_1]
			+
			\cdots
			+
			E[a_nX_n]
			\\[0.3cm]
			&=
			a_1E[X_1]
			+\cdots+
			a_nE[X_n].
		\end{aligned}
		\]
		
		Hence,
		
		\[
		\boxed{
			E\!\left[
			\sum_{i=1}^{n}
			a_iX_i
			\right]
			=
			\sum_{i=1}^{n}
			a_iE[X_i].
		}
		\]
		
		\bigskip
		
		\paragraph{Special Cases}
		
		Several useful identities are immediate consequences of linearity.
		
		Choosing
		
		\[
		a=1,
		\qquad
		b=1,
		\]
		
		gives
		
		\[
		\boxed{
			E[X+Y]
			=
			E[X]
			+
			E[Y].
		}
		\]
		
		Choosing
		
		\[
		a=1,
		\qquad
		b=-1,
		\]
		
		gives
		
		\[
		\boxed{
			E[X-Y]
			=
			E[X]
			-
			E[Y].
		}
		\]
		
		Choosing
		
		\[
		b=0,
		\]
		
		gives
		
		\[
		\boxed{
			E[aX]
			=
			aE[X].
		}
		\]
		
		Choosing
		
		\[
		a=1,
		\]
		
		and letting \(Y=c\) be a constant,
		
		\[
		\begin{aligned}
			E[X+c]
			&=
			E[X]
			+
			E[c]
			\\[0.3cm]
			&=
			E[X]
			+
			c.
		\end{aligned}
		\]
		
		Thus,
		
		\[
		\boxed{
			E[X+c]
			=
			E[X]
			+
			c.
		}
		\]
		
		\bigskip
		
		\paragraph{Example}
		
		Suppose
		
		\[
		E[X]=5,
		\qquad
		E[Y]=2.
		\]
		
		Compute
		
		\[
		E[3X-4Y+7].
		\]
		
		Applying linearity,
		
		\[
		\begin{aligned}
			E[3X-4Y+7]
			&=
			3E[X]
			-
			4E[Y]
			+
			E[7]
			\\[0.3cm]
			&=
			3(5)
			-
			4(2)
			+
			7
			\\[0.3cm]
			&=
			15
			-
			8
			+
			7
			\\[0.3cm]
			&=
			14.
		\end{aligned}
		\]
		
		\[
		\boxed{
			E[3X-4Y+7]=14.
		}
		\]
		
		\bigskip
		
		\paragraph{Importance of Linearity}
		
		Linearity of expectation is one of the most powerful tools in probability
		theory because it holds without requiring independence among the random
		variables.
		
		This property greatly simplifies the computation of expectations and plays
		a central role in statistics.
		
		For example, it allows us to show that the sample mean,
		
		\[
		\bar X
		=
		\frac1n
		\sum_{i=1}^{n}
		X_i,
		\]
		
		satisfies
		
		\[
		E[\bar X]
		=
		\frac1n
		\sum_{i=1}^{n}
		E[X_i].
		\]
		
		If the observations are identically distributed with mean
		
		\[
		\mu,
		\]
		
		then
		
		\[
		E[\bar X]
		=
		\frac1n(n\mu)
		=
		\mu,
		\]
		
		showing that the sample mean is an unbiased estimator of the population
		mean.
		
	\end{mdframed}
	
	\section{Expectation of Sums}
	\begin{mdframed}[
		linewidth=0.8pt,
		roundcorner=6pt,
		backgroundcolor=gray!5,
		linecolor=gray!60,
		innertopmargin=10pt,
		innerbottommargin=10pt,
		innerleftmargin=10pt,
		innerrightmargin=10pt
		]
		
		\subsection*{Expectation of Sums}
		
		One of the most useful consequences of the linearity of expectation is that
		the expectation of the sum of several random variables is equal to the sum
		of their expectations.
		
		This property is extremely important in probability and statistics because
		it allows us to compute the expected value of complicated random quantities
		by studying each component separately.
		
		Unlike many other probabilistic results, \emph{no independence assumptions
			are required}. The result is true for every collection of random variables
		whose expectations exist.
		
		\bigskip
		
		\paragraph{Theorem}
		
		Let
		
		\[
		X_1,X_2,\ldots,X_n
		\]
		
		be random variables with finite expectations.
		
		Then
		
		\[
		\boxed{
			E[X_1+X_2+\cdots+X_n]
			=
			E[X_1]+E[X_2]+\cdots+E[X_n].
		}
		\]
		
		Equivalently,
		
		\[
		\boxed{
			E\!\left(
			\sum_{i=1}^{n}X_i
			\right)
			=
			\sum_{i=1}^{n}E[X_i].
		}
		\]
		
		\bigskip
		
		\paragraph{Proof}
		
		The result follows directly from the linearity of expectation.
		
		Beginning with
		
		\[
		X_1+X_2+\cdots+X_n,
		\]
		
		we apply the linearity property repeatedly.
		
		For two random variables,
		
		\[
		E[X_1+X_2]
		=
		E[X_1]
		+
		E[X_2].
		\]
		
		Applying linearity once more,
		
		\[
		\begin{aligned}
			E[X_1+X_2+X_3]
			&=
			E[(X_1+X_2)+X_3]
			\\[0.3cm]
			&=
			E[X_1+X_2]
			+
			E[X_3]
			\\[0.3cm]
			&=
			E[X_1]
			+
			E[X_2]
			+
			E[X_3].
		\end{aligned}
		\]
		
		Continuing this process,
		
		\[
		\begin{aligned}
			E[X_1+\cdots+X_n]
			&=
			E[(X_1+\cdots+X_{n-1})+X_n]
			\\[0.3cm]
			&=
			E[X_1+\cdots+X_{n-1}]
			+
			E[X_n]
			\\[0.3cm]
			&=
			E[X_1]
			+\cdots+
			E[X_n].
		\end{aligned}
		\]
		
		Hence,
		
		\[
		\boxed{
			E\!\left(
			\sum_{i=1}^{n}X_i
			\right)
			=
			\sum_{i=1}^{n}E[X_i].
		}
		\]
		
		\bigskip
		
		\paragraph{An Alternative Proof}
		
		The result also follows immediately from the general linearity theorem.
		
		Indeed,
		
		\[
		\sum_{i=1}^{n}X_i
		=
		1X_1
		+
		1X_2
		+\cdots+
		1X_n.
		\]
		
		Applying
		
		\[
		E\!\left(
		\sum_{i=1}^{n}
		a_iX_i
		\right)
		=
		\sum_{i=1}^{n}
		a_iE[X_i],
		\]
		
		with
		
		\[
		a_i=1
		\qquad
		(i=1,\ldots,n),
		\]
		
		gives
		
		\[
		E\!\left(
		\sum_{i=1}^{n}X_i
		\right)
		=
		\sum_{i=1}^{n}E[X_i].
		\]
		
		\bigskip
		
		\paragraph{No Independence Is Required}
		
		A common misconception is that the random variables must be independent for
		linearity of expectation to hold.
		
		This is \textbf{not true}.
		
		The identity
		
		\[
		\boxed{
			E[X_1+\cdots+X_n]
			=
			E[X_1]+\cdots+E[X_n]
		}
		\]
		
		holds for \emph{every} collection of random variables whose expectations
		exist, regardless of whether they are independent, correlated, or even
		deterministically related.
		
		The reason is that expectation is a \emph{linear operator}. Its definition
		depends only on averaging the values of the random variables and does not
		involve any interaction terms between them.
		
		\bigskip
		
		To illustrate this, suppose that
		
		\[
		X_1=X_2=\cdots=X_n=X.
		\]
		
		The random variables are perfectly positively correlated; knowing one of
		them determines all the others exactly.
		
		Then,
		
		\[
		X_1+\cdots+X_n
		=
		nX.
		\]
		
		Applying the linearity of expectation,
		
		\[
		\begin{aligned}
			E[X_1+\cdots+X_n]
			&=
			E[nX]
			\\[0.3cm]
			&=
			nE[X].
		\end{aligned}
		\]
		
		On the other hand,
		
		\[
		\begin{aligned}
			E[X_1]+\cdots+E[X_n]
			&=
			E[X]+\cdots+E[X]
			\\[0.3cm]
			&=
			nE[X].
		\end{aligned}
		\]
		
		Hence,
		
		\[
		\boxed{
			E[X_1+\cdots+X_n]
			=
			E[X_1]+\cdots+E[X_n],
		}
		\]
		
		even though the variables are as dependent as possible.
		
		\bigskip
		
		\paragraph{What Changes When Variables Are Dependent?}
		
		Although dependence does not affect the expectation of a sum, it does affect
		its variance.
		
		Recall that
		
		\[
		\operatorname{Var}(X)
		=
		E[(X-E[X])^2].
		\]
		
		For two random variables,
		
		\[
		\boxed{
			\operatorname{Var}(X+Y)
			=
			\operatorname{Var}(X)
			+
			\operatorname{Var}(Y)
			+
			2\,\operatorname{Cov}(X,Y),
		}
		\]
		
		where
		
		\[
		\operatorname{Cov}(X,Y)
		=
		E[(X-E[X])(Y-E[Y])]
		\]
		
		is the covariance between \(X\) and \(Y\).
		
		linewidth=0.8pt,
		roundcorner=6pt,
		backgroundcolor=gray!5,
		linecolor=gray!60,
		innertopmargin=10pt,
		innerbottommargin=10pt,
		innerleftmargin=10pt,
		innerrightmargin=10pt
		]
		
		\paragraph{Proof of the Variance of a Sum}
		
		We wish to prove that
		
		\[
		\boxed{
			\operatorname{Var}(X+Y)
			=
			\operatorname{Var}(X)
			+
			\operatorname{Var}(Y)
			+
			2\,\operatorname{Cov}(X,Y).
		}
		\]
		
		This identity is valid for \emph{all} random variables having finite second
		moments. No independence assumption is required.
		
		\bigskip
		
		\paragraph{Step 1: Recall the Definition of Variance}
		
		By definition,
		
		\[
		\operatorname{Var}(Z)
		=
		E\!\left[(Z-E[Z])^2\right].
		\]
		
		Applying this definition to the random variable
		
		\[
		Z=X+Y,
		\]
		
		we obtain
		
		\[
		\operatorname{Var}(X+Y)
		=
		E\!\left[
		\left(
		(X+Y)-E[X+Y]
		\right)^2
		\right].
		\]
		
		\bigskip
		
		\paragraph{Step 2: Use Linearity of Expectation}
		
		Since expectation is linear,
		
		\[
		E[X+Y]
		=
		E[X]
		+
		E[Y].
		\]
		
		Substituting,
		
		\[
		\operatorname{Var}(X+Y)
		=
		E\!\left[
		\left(
		X+Y-E[X]-E[Y]
		\right)^2
		\right].
		\]
		
		Grouping terms,
		
		\[
		\operatorname{Var}(X+Y)
		=
		E\!\left[
		\left(
		(X-E[X])
		+
		(Y-E[Y])
		\right)^2
		\right].
		\]
		
		\bigskip
		
		\paragraph{Step 3: Expand the Square}
		
		Recall the algebraic identity
		
		\[
		(a+b)^2
		=
		a^2+2ab+b^2.
		\]
		
		Let
		
		\[
		a=X-E[X],
		\qquad
		b=Y-E[Y].
		\]
		
		Then
		
		\[
		\begin{aligned}
			\operatorname{Var}(X+Y)
			&=
			E\!\left[
			(X-E[X])^2
			+
			2(X-E[X])(Y-E[Y])
			+
			(Y-E[Y])^2
			\right].
		\end{aligned}
		\]
		
		\bigskip
		
		\paragraph{Step 4: Apply Linearity of Expectation}
		
		Using the linearity of expectation,
		
		\[
		\begin{aligned}
			\operatorname{Var}(X+Y)
			&=
			E[(X-E[X])^2]
			\\
			&\quad
			+
			2E[(X-E[X])(Y-E[Y])]
			\\
			&\quad
			+
			E[(Y-E[Y])^2].
		\end{aligned}
		\]
		
		\bigskip
		
		\paragraph{Step 5: Recognize Each Term}
		
		By definition,
		
		\[
		E[(X-E[X])^2]
		=
		\operatorname{Var}(X),
		\]
		
		and
		
		\[
		E[(Y-E[Y])^2]
		=
		\operatorname{Var}(Y).
		\]
		
		Moreover, the covariance is defined as
		
		\[
		\boxed{
			\operatorname{Cov}(X,Y)
			=
			E[(X-E[X])(Y-E[Y])].
		}
		\]
		
		Therefore,
		
		\[
		E[(X-E[X])(Y-E[Y])]
		=
		\operatorname{Cov}(X,Y).
		\]
		
		Substituting these identities,
		
		\[
		\begin{aligned}
			\operatorname{Var}(X+Y)
			&=
			\operatorname{Var}(X)
			+
			2\operatorname{Cov}(X,Y)
			+
			\operatorname{Var}(Y).
		\end{aligned}
		\]
		
		Rearranging the terms,
		
		\[
		\boxed{
			\operatorname{Var}(X+Y)
			=
			\operatorname{Var}(X)
			+
			\operatorname{Var}(Y)
			+
			2\,\operatorname{Cov}(X,Y).
		}
		\]
		
		This completes the proof.
		
		\bigskip
		
		\paragraph{Special Case: Independent Random Variables}
		
		If \(X\) and \(Y\) are independent, then
		
		\[
		\operatorname{Cov}(X,Y)=0.
		\]
		
		Consequently,
		
		\[
		\boxed{
			\operatorname{Var}(X+Y)
			=
			\operatorname{Var}(X)
			+
			\operatorname{Var}(Y).
		}
		\]
		
		This is the familiar variance formula for independent random variables.
		
		\bigskip
		
		\paragraph{Remark}
		
		This proof highlights an important contrast between expectation and variance.
		
		Expectation is always linear:
		
		\[
		E[X+Y]
		=
		E[X]
		+
		E[Y],
		\]
		
		regardless of any dependence between the variables.
		
		Variance, however, is generally \emph{not} linear because of the covariance term, $	2\,\operatorname{Cov}(X,Y),	$ which measures the extent to which the two random variables vary together.
		Only when this covariance is zero (for example, under independence) does the variance of a sum reduce to the sum of the variances.
		
		More generally, for \(n\) random variables,
		
		\[
		\boxed{
			\operatorname{Var}
			\!\left(
			\sum_{i=1}^{n}X_i
			\right)
			=
			\sum_{i=1}^{n}
			\operatorname{Var}(X_i)
			+
			2
			\sum_{i<j}
			\operatorname{Cov}(X_i,X_j).
		}
		\]
		
		Unlike expectation, the variance depends explicitly on the covariance
		between the variables.
		
		\bigskip
		
		\paragraph{The Independent Case}
		
		If the random variables are independent, then
		
		\[
		\operatorname{Cov}(X_i,X_j)=0,
		\qquad
		i\neq j.
		\]
		
		Consequently,
		
		\[
		\boxed{
			\operatorname{Var}
			\!\left(
			\sum_{i=1}^{n}X_i
			\right)
			=
			\sum_{i=1}^{n}
			\operatorname{Var}(X_i).
		}
		\]
		
		This simplification is one of the main reasons why independence is such an
		important assumption in probability and statistics.
		
		\bigskip
		
		\paragraph{The Dependent Case}
		
		When the variables are dependent, the covariance terms generally do not
		vanish.
		
		For example, suppose
		
		\[
		X_1=X_2=\cdots=X_n=X.
		\]
		
		Then
		
		\[
		\sum_{i=1}^{n}X_i=nX,
		\]
		
		so
		
		\[
		\begin{aligned}
			\operatorname{Var}(X_1+\cdots+X_n)
			&=
			\operatorname{Var}(nX)
			\\[0.3cm]
			&=
			n^2\operatorname{Var}(X).
		\end{aligned}
		\]
		
		However,
		
		\[
		\sum_{i=1}^{n}
		\operatorname{Var}(X_i)
		=
		n\operatorname{Var}(X),
		\]
		
		which is much smaller whenever \(n>1\).
		
		The difference is explained by the covariance terms.
		
		Indeed,
		
		\[
		\operatorname{Cov}(X_i,X_j)
		=
		\operatorname{Var}(X),
		\qquad
		i\neq j,
		\]
		
		so
		
		\[
		\begin{aligned}
			\operatorname{Var}
			\!\left(
			\sum_{i=1}^{n}X_i
			\right)
			&=
			n\operatorname{Var}(X)
			+
			2
			\binom{n}{2}
			\operatorname{Var}(X)
			\\[0.3cm]
			&=
			n\operatorname{Var}(X)
			+
			n(n-1)\operatorname{Var}(X)
			\\[0.3cm]
			&=
			n^2\operatorname{Var}(X),
		\end{aligned}
		\]
		
		which agrees with the direct calculation.
		
		\bigskip
		
		\paragraph{Conclusion}
		
		The expectation operator is always linear:
		
		\[
		\boxed{
			E\!\left[
			\sum_{i=1}^{n}X_i
			\right]
			=
			\sum_{i=1}^{n}E[X_i],
		}
		\]
		
		regardless of the dependence among the random variables.
		
		In contrast, the variance of a sum depends on the covariance structure:
		
		\[
		\boxed{
			\operatorname{Var}
			\!\left(
			\sum_{i=1}^{n}X_i
			\right)
			=
			\sum_{i=1}^{n}\operatorname{Var}(X_i)
			+
			2\sum_{i<j}\operatorname{Cov}(X_i,X_j).
		}
		\]
		
		Thus, independence is **not** required for linearity of expectation, but it
		is essential for the simple variance formula
		
		\[
		\boxed{
			\operatorname{Var}
			\!\left(
			\sum_{i=1}^{n}X_i
			\right)
			=
			\sum_{i=1}^{n}\operatorname{Var}(X_i).
		}
		\]
		
		This distinction is one of the most important conceptual differences between
		expectation and variance.
		
		\bigskip
		
		\paragraph{Example 1}
		
		Suppose
		
		\[
		E[X_1]=3,
		\qquad
		E[X_2]=-1,
		\qquad
		E[X_3]=5.
		\]
		
		Then
		
		\[
		\begin{aligned}
			E[X_1+X_2+X_3]
			&=
			3
			+
			(-1)
			+
			5
			\\[0.3cm]
			&=
			7.
		\end{aligned}
		\]
		
		Therefore,
		
		\[
		\boxed{
			E[X_1+X_2+X_3]=7.
		}
		\]
		
		\bigskip
		
		\paragraph{Example 2}
		
		Suppose
		
		\[
		E[X_i]=10,
		\qquad
		i=1,\ldots,100.
		\]
		
		Then
		
		\[
		\begin{aligned}
			E[X_1+\cdots+X_{100}]
			&=
			100\times10
			\\[0.3cm]
			&=
			1000.
		\end{aligned}
		\]
		
		Notice that we never needed to know whether the variables were
		independent.
		
		\bigskip
		
		\paragraph{Application to the Sample Mean}
		
		One of the most important applications of this theorem is the computation
		of the expected value of the sample mean.
		
		Recall that
		
		\[
		\bar X
		=
		\frac1n
		\sum_{i=1}^{n}X_i.
		\]
		
		Applying the linearity of expectation,
		
		\[
		\begin{aligned}
			E[\bar X]
			&=
			E\!\left(
			\frac1n
			\sum_{i=1}^{n}X_i
			\right)
			\\[0.3cm]
			&=
			\frac1n
			E\!\left(
			\sum_{i=1}^{n}X_i
			\right)
			\\[0.3cm]
			&=
			\frac1n
			\sum_{i=1}^{n}
			E[X_i].
		\end{aligned}
		\]
		
		If the observations are identically distributed with common mean
		
		\[
		E[X_i]=\mu,
		\]
		
		then
		
		\[
		\begin{aligned}
			E[\bar X]
			&=
			\frac1n
			(\mu+\mu+\cdots+\mu)
			\\[0.3cm]
			&=
			\frac1n(n\mu)
			\\[0.3cm]
			&=
			\mu.
		\end{aligned}
		\]
		
		Therefore,
		
		\[
		\boxed{
			E[\bar X]=\mu.
		}
		\]
		
		This shows that the sample mean is an unbiased estimator of the population
		mean.
		
		\bigskip
		
		\paragraph{Importance of the Result}
		
		The expectation of sums is one of the most frequently used tools in
		probability and statistics.
		
		It appears in the derivation of the sample mean, estimation theory, the Law
		of Large Numbers, the Central Limit Theorem, regression analysis, stochastic
		processes, and many other areas.
		
		Its simplicity is one of the reasons why expectation is often much easier
		to analyze than other quantities such as the variance, whose behavior under
		addition requires additional assumptions about covariance and
		independence.
		
	\end{mdframed}
	
	\section{Expectation of Products of Independent Variables}
	\begin{mdframed}[
		linewidth=0.8pt,
		roundcorner=6pt,
		backgroundcolor=gray!5,
		linecolor=gray!60,
		innertopmargin=10pt,
		innerbottommargin=10pt,
		innerleftmargin=10pt,
		innerrightmargin=10pt
		]
		
		\subsection*{Expectation of Products of Independent Random Variables}
		
		Unlike the expectation of a sum, the expectation of a product generally
		cannot be separated into the product of the expectations.
		
		In fact,
		
		\[
		E[XY]
		\neq
		E[X]E[Y]
		\]
		
		in general.
		
		However, if the random variables are independent, then a remarkable
		simplification occurs.
		
		\bigskip
		
		\paragraph{Theorem}
		
		Let
		
		\[
		X
		\quad\text{and}\quad
		Y
		\]
		
		be independent random variables whose expectations exist.
		
		Then
		
		\[
		\boxed{
			E[XY]
			=
			E[X]E[Y].
		}
		\]
		
		More generally, if
		
		\[
		X_1,X_2,\ldots,X_n
		\]
		
		are mutually independent random variables with finite expectations, then
		
		\[
		\boxed{
			E\!\left(
			\prod_{i=1}^{n}X_i
			\right)
			=
			\prod_{i=1}^{n}E[X_i].
		}
		\]
		
		\bigskip
		
		\paragraph{Proof for Discrete Random Variables}
		
		Suppose that
		
		\[
		X
		\]
		
		and
		
		\[
		Y
		\]
		
		are discrete random variables.
		
		By the definition of expectation,
		
		\[
		E[XY]
		=
		\sum_x
		\sum_y
		xyP(X=x,Y=y).
		\]
		
		Since \(X\) and \(Y\) are independent,
		
		\[
		P(X=x,Y=y)
		=
		P(X=x)P(Y=y).
		\]
		
		Substituting this into the expectation,
		
		\[
		\begin{aligned}
			E[XY]
			&=
			\sum_x
			\sum_y
			xyP(X=x)P(Y=y).
		\end{aligned}
		\]
		
		Rearranging the terms,
		
		\[
		\begin{aligned}
			E[XY]
			&=
			\sum_x
			xP(X=x)
			\sum_y
			yP(Y=y).
		\end{aligned}
		\]
		
		The first summation is simply
		
		\[
		E[X],
		\]
		
		while the second summation is
		
		\[
		E[Y].
		\]
		
		Therefore,
		
		\[
		\boxed{
			E[XY]
			=
			E[X]E[Y].
		}
		\]
		
		\bigskip
		
		\paragraph{Proof for Continuous Random Variables}
		
		Suppose now that
		
		\[
		X
		\]
		
		and
		
		\[
		Y
		\]
		
		are continuous random variables.
		
		The expectation of their product is
		
		\[
		E[XY]
		=
		\int_{-\infty}^{\infty}
		\int_{-\infty}^{\infty}
		xy
		f_{X,Y}(x,y)
		\,dx\,dy.
		\]
		
		Since \(X\) and \(Y\) are independent,
		
		\[
		f_{X,Y}(x,y)
		=
		f_X(x)f_Y(y).
		\]
		
		Substituting,
		
		\[
		\begin{aligned}
			E[XY]
			&=
			\int_{-\infty}^{\infty}
			\int_{-\infty}^{\infty}
			xy
			f_X(x)
			f_Y(y)
			\,dx\,dy.
		\end{aligned}
		\]
		
		Since the integrand factors into a function of \(x\) multiplied by a
		function of \(y\), the double integral separates into the product of two
		single integrals,
		
		\[
		\begin{aligned}
			E[XY]
			&=
			\left(
			\int_{-\infty}^{\infty}
			xf_X(x)\,dx
			\right)
			\left(
			\int_{-\infty}^{\infty}
			yf_Y(y)\,dy
			\right).
		\end{aligned}
		\]
		
		Recognizing these integrals,
		
		\[
		\boxed{
			E[XY]
			=
			E[X]E[Y].
		}
		\]
		
		\bigskip
		
		\paragraph{Generalization to Several Random Variables}
		
		Suppose
		
		\[
		X_1,X_2,\ldots,X_n
		\]
		
		are mutually independent.
		
		Then
		
		\[
		E[X_1X_2\cdots X_n]
		=
		\sum\cdots\sum
		x_1x_2\cdots x_n
		P(X_1=x_1,\ldots,X_n=x_n).
		\]
		
		Independence implies
		
		\[
		P(X_1=x_1,\ldots,X_n=x_n)
		=
		\prod_{i=1}^{n}
		P(X_i=x_i).
		\]
		
		Hence,
		
		\[
		\begin{aligned}
			E[X_1X_2\cdots X_n]
			&=
			\sum\cdots\sum
			\left(
			x_1P(X_1=x_1)
			\right)
			\cdots
			\left(
			x_nP(X_n=x_n)
			\right).
		\end{aligned}
		\]
		
		The multiple summation separates into the product of the individual
		summations,
		
		\[
		\begin{aligned}
			E[X_1X_2\cdots X_n]
			&=
			\left(
			\sum x_1P(X_1=x_1)
			\right)
			\cdots
			\left(
			\sum x_nP(X_n=x_n)
			\right)
			\\[0.3cm]
			&=
			E[X_1]
			E[X_2]
			\cdots
			E[X_n].
		\end{aligned}
		\]
		
		Therefore,
		
		\[
		\boxed{
			E\!\left(
			\prod_{i=1}^{n}X_i
			\right)
			=
			\prod_{i=1}^{n}E[X_i].
		}
		\]
		
		\bigskip
		
		\paragraph{Why Independence Is Necessary}
		
		The previous theorem depends crucially on the independence of the random
		variables.
		
		If the variables are dependent, the identity generally fails.
		
		For example, let
		
		\[
		Y=X.
		\]
		
		Then
		
		\[
		XY=X^2,
		\]
		
		so that
		
		\[
		E[XY]
		=
		E[X^2].
		\]
		
		On the other hand,
		
		\[
		E[X]E[Y]
		=
		(E[X])^2.
		\]
		
		These two quantities are generally different because
		
		\[
		E[X^2]
		=
		(E[X])^2+\operatorname{Var}(X).
		\]
		
		Thus,
		
		\[
		E[X^2]
		=
		(E[X])^2
		\]
		
		only when
		
		\[
		\operatorname{Var}(X)=0,
		\]
		
		that is, when \(X\) is a constant random variable.
		
		This example illustrates that independence is essential.
		
		\bigskip
		
		\paragraph{Example}
		
		Suppose
		
		\[
		E[X]=4,
		\qquad
		E[Y]=3,
		\]
		
		and assume that \(X\) and \(Y\) are independent.
		
		Then
		
		\[
		\begin{aligned}
			E[XY]
			&=
			E[X]E[Y]
			\\[0.3cm]
			&=
			4\times3
			\\[0.3cm]
			&=
			12.
		\end{aligned}
		\]
		
		Hence,
		
		\[
		\boxed{
			E[XY]=12.
		}
		\]
		
		\bigskip
		
		\paragraph{Application to Variance}
		
		One of the most important applications of this theorem occurs in the study
		of variance.
		
		Recall that
		
		\[
		\operatorname{Var}(X)
		=
		E[X^2]-(E[X])^2.
		\]
		
		If \(X\) and \(Y\) are independent,
		
		\[
		E[XY]
		=
		E[X]E[Y].
		\]
		
		Consequently,
		
		\[
		\begin{aligned}
			\operatorname{Var}(X+Y)
			&=
			E[(X+Y)^2]
			-
			(E[X+Y])^2
			\\[0.3cm]
			&=
			E[X^2]
			+
			2E[XY]
			+
			E[Y^2]
			\\
			&\qquad
			-
			(E[X]+E[Y])^2.
		\end{aligned}
		\]
		
		Using
		
		\[
		E[XY]=E[X]E[Y],
		\]
		
		the cross terms cancel,
		
		\[
		\boxed{
			\operatorname{Var}(X+Y)
			=
			\operatorname{Var}(X)
			+
			\operatorname{Var}(Y).
		}
		\]
		
		This important identity would generally be false without the independence
		assumption.
		
		\bigskip
		
		\paragraph{Summary}
		
		For sums,
		
		\[
		\boxed{
			E[X+Y]
			=
			E[X]
			+
			E[Y],
		}
		\]
		
		and this holds for all random variables.
		
		For products,
		
		\[
		\boxed{
			E[XY]
			=
			E[X]E[Y],
		}
		\]
		
		the identity holds only when the random variables are independent.
		
		This distinction is one of the most important facts in probability theory
		and should always be kept in mind when manipulating expected values.
		
	\end{mdframed}
	
	
	\section{Monotonicity}
	\begin{mdframed}[
		linewidth=0.8pt,
		roundcorner=6pt,
		backgroundcolor=gray!5,
		linecolor=gray!60,
		innertopmargin=10pt,
		innerbottommargin=10pt,
		innerleftmargin=10pt,
		innerrightmargin=10pt
		]
		
		\subsection*{Monotonicity of Expectation}
		
		One of the most natural properties of the expected value is that it preserves
		order.
		
		Intuitively, if one random variable is always less than or equal to another,
		then its average value should also be less than or equal to the average of
		the larger variable.
		
		This property is known as the \emph{monotonicity of expectation}.
		
		\bigskip
		
		\paragraph{Theorem (Monotonicity)}
		
		Let
		
		\[
		X
		\quad\text{and}\quad
		Y
		\]
		
		be random variables such that
		
		\[
		X\le Y
		\]
		
		with probability one.
		
		Then
		
		\[
		\boxed{
			E[X]
			\le
			E[Y].
		}
		\]
		
		Equivalently,
		
		\[
		X\le Y
		\quad\Longrightarrow\quad
		E[X]\le E[Y].
		\]
		
		\bigskip
		
		\paragraph{Interpretation}
		
		The condition
		
		\[
		X\le Y
		\]
		
		means that for every possible outcome,
		
		\[
		X(\omega)\le Y(\omega).
		\]
		
		That is, no realization of \(X\) ever exceeds the corresponding realization
		of \(Y\).
		
		Since expectation is an average over all possible realizations, it is
		natural that the average value of \(X\) cannot exceed the average value of
		\(Y\).
		
		\bigskip
		
		\paragraph{Proof Using Linearity}
		
		Define the random variable
		
		\[
		Z=Y-X.
		\]
		
		Since
		
		\[
		X\le Y,
		\]
		
		we have
		
		\[
		Z\ge0
		\]
		
		for every outcome.
		
		Applying the linearity of expectation,
		
		\[
		\begin{aligned}
			E[Z]
			&=
			E[Y-X]
			\\[0.3cm]
			&=
			E[Y]-E[X].
		\end{aligned}
		\]
		
		Since \(Z\) is always nonnegative,
		
		\[
		E[Z]\ge0.
		\]
		
		Therefore,
		
		\[
		E[Y]-E[X]\ge0,
		\]
		
		which implies
		
		\[
		\boxed{
			E[X]\le E[Y].
		}
		\]
		
		This completes the proof.
		
		\bigskip
		
		\paragraph{Direct Proof for Discrete Random Variables}
		
		Suppose that
		
		\[
		X
		\]
		
		has probability mass function
		
		\[
		P(X=x_i)=p_i,
		\]
		
		and that
		
		\[
		g(X)=Y,
		\]
		
		where
		
		\[
		g(x)\ge x
		\]
		
		for every possible value of \(x\).
		
		Then
		
		\[
		\begin{aligned}
			E[Y]-E[X]
			&=
			\sum_i
			g(x_i)p_i
			-
			\sum_i
			x_ip_i
			\\[0.3cm]
			&=
			\sum_i
			(g(x_i)-x_i)p_i.
		\end{aligned}
		\]
		
		Since
		
		\[
		g(x_i)-x_i\ge0,
		\]
		
		and
		
		\[
		p_i\ge0,
		\]
		
		every term in the summation is nonnegative.
		
		Hence,
		
		\[
		E[Y]-E[X]\ge0,
		\]
		
		or equivalently,
		
		\[
		\boxed{
			E[X]\le E[Y].
		}
		\]
		
		\bigskip
		
		\paragraph{Direct Proof for Continuous Random Variables}
		
		Suppose
		
		\[
		X\le Y
		\]
		
		almost surely.
		
		Then
		
		\[
		Y-X\ge0.
		\]
		
		Therefore,
		
		\[
		\begin{aligned}
			E[Y]-E[X]
			&=
			E[Y-X]
			\\[0.3cm]
			&=
			\int_{-\infty}^{\infty}
			(Y-X)f(x)\,dx.
		\end{aligned}
		\]
		
		Since
		
		\[
		Y-X\ge0
		\]
		
		and
		
		\[
		f(x)\ge0,
		\]
		
		the integrand is everywhere nonnegative.
		
		Consequently,
		
		\[
		E[Y]-E[X]\ge0,
		\]
		
		which again yields
		
		\[
		\boxed{
			E[X]\le E[Y].
		}
		\]
		
		\bigskip
		
		\paragraph{Example 1}
		
		Suppose
		
		\[
		X=
		\begin{cases}
			1,&P=\dfrac12,\\
			3,&P=\dfrac12,
		\end{cases}
		\]
		
		and define
		
		\[
		Y=X+2.
		\]
		
		Clearly,
		
		\[
		X\le Y
		\]
		
		for every outcome.
		
		Now compute the expectations.
		
		First,
		
		\[
		\begin{aligned}
			E[X]
			&=
			1\left(\frac12\right)
			+
			3\left(\frac12\right)
			\\[0.3cm]
			&=
			2.
		\end{aligned}
		\]
		
		Next,
		
		\[
		\begin{aligned}
			E[Y]
			&=
			E[X+2]
			\\[0.3cm]
			&=
			E[X]+2
			\\[0.3cm]
			&=
			4.
		\end{aligned}
		\]
		
		Thus,
		
		\[
		\boxed{
			E[X]=2<4=E[Y].
		}
		\]
		
		\bigskip
		
		\paragraph{Example 2}
		
		Suppose
		
		\[
		0\le X\le10.
		\]
		
		Applying monotonicity twice,
		
		\[
		E[0]
		\le
		E[X]
		\le
		E[10].
		\]
		
		Since
		
		\[
		E[0]=0,
		\qquad
		E[10]=10,
		\]
		
		we conclude
		
		\[
		\boxed{
			0
			\le
			E[X]
			\le
			10.
		}
		\]
		
		Thus, whenever a random variable is bounded, its expectation must lie
		between the same bounds.
		
		\bigskip
		
		\paragraph{Important Consequences}
		
		Monotonicity immediately implies several useful inequalities.
		
		If
		
		\[
		X\ge0,
		\]
		
		then
		
		\[
		0\le X.
		\]
		
		Applying monotonicity,
		
		\[
		\boxed{
			E[X]\ge0.
		}
		\]
		
		Similarly, if
		
		\[
		a\le X\le b,
		\]
		
		then
		
		\[
		\boxed{
			a
			\le
			E[X]
			\le
			b.
		}
		\]
		
		These results are frequently used when proving bounds in probability and
		statistics.
		
		\bigskip
		
		\paragraph{Importance of Monotonicity}
		
		Monotonicity is one of the fundamental properties of expectation.
		
		Together with linearity, it forms the basis of many important inequalities,
		including
		
		\begin{itemize}
			\item Markov's Inequality,
			\item Chebyshev's Inequality,
			\item Jensen's Inequality,
			\item Hölder's Inequality,
			\item Minkowski's Inequality.
		\end{itemize}
		
		Many advanced results in probability theory ultimately rely on these two
		simple properties of expectation.
		
	\end{mdframed}
	
	
	\section{Non-negativity}
	\begin{mdframed}[
		linewidth=0.8pt,
		roundcorner=6pt,
		backgroundcolor=gray!5,
		linecolor=gray!60,
		innertopmargin=10pt,
		innerbottommargin=10pt,
		innerleftmargin=10pt,
		innerrightmargin=10pt
		]
		
		\subsection*{Non-negativity of Expectation}
		
		An immediate consequence of the monotonicity property is that the expected
		value of a nonnegative random variable is itself nonnegative.
		
		This result is one of the most frequently used properties of expectation,
		particularly in proving inequalities and establishing bounds in probability
		theory.
		
		\bigskip
		
		\paragraph{Theorem (Non-negativity)}
		
		Let
		
		\[
		X
		\]
		
		be a random variable satisfying
		
		\[
		X\ge0
		\]
		
		almost surely.
		
		Then
		
		\[
		\boxed{
			E[X]\ge0.
		}
		\]
		
		\bigskip
		
		\paragraph{Interpretation}
		
		The condition
		
		\[
		X\ge0
		\]
		
		means that every possible realization of the random variable is
		nonnegative.
		
		Since expectation is simply a weighted average of all possible values, an
		average of nonnegative numbers can never be negative.
		
		Therefore,
		
		\[
		E[X]\ge0.
		\]
		
		\bigskip
		
		\paragraph{Proof Using Monotonicity}
		
		Since
		
		\[
		0\le X,
		\]
		
		the monotonicity of expectation immediately gives
		
		\[
		E[0]
		\le
		E[X].
		\]
		
		From the property of the expectation of a constant,
		
		\[
		E[0]=0.
		\]
		
		Hence,
		
		\[
		\boxed{
			E[X]\ge0.
		}
		\]
		
		This proof is valid for every random variable whose expectation exists.
		
		\bigskip
		
		\paragraph{Direct Proof for Discrete Random Variables}
		
		Suppose that
		
		\[
		X
		\]
		
		is discrete with probability mass function
		
		\[
		P(X=x_i)=p_i.
		\]
		
		Since
		
		\[
		X\ge0,
		\]
		
		every possible value satisfies
		
		\[
		x_i\ge0.
		\]
		
		Moreover,
		
		\[
		p_i\ge0
		\]
		
		for every \(i\).
		
		The expectation is
		
		\[
		E[X]
		=
		\sum_i
		x_ip_i.
		\]
		
		Each term
		
		\[
		x_ip_i
		\]
		
		is the product of two nonnegative numbers.
		
		Therefore,
		
		\[
		x_ip_i\ge0
		\]
		
		for every \(i\).
		
		Since a sum of nonnegative numbers is nonnegative,
		
		\[
		\boxed{
			E[X]\ge0.
		}
		\]
		
		\bigskip
		
		\paragraph{Direct Proof for Continuous Random Variables}
		
		Suppose that
		
		\[
		X
		\]
		
		has probability density function
		
		\[
		f_X(x).
		\]
		
		The expectation is
		
		\[
		E[X]
		=
		\int_{-\infty}^{\infty}
		x
		f_X(x)
		\,dx.
		\]
		
		Since
		
		\[
		X\ge0
		\]
		
		almost surely,
		
		\[
		f_X(x)=0
		\qquad
		\text{whenever }x<0.
		\]
		
		Therefore,
		
		\[
		E[X]
		=
		\int_0^{\infty}
		x
		f_X(x)
		\,dx.
		\]
		
		Both factors satisfy
		
		\[
		x\ge0,
		\qquad
		f_X(x)\ge0.
		\]
		
		Hence,
		
		\[
		x
		f_X(x)
		\ge0
		\]
		
		for every \(x\).
		
		The integral of a nonnegative function is itself nonnegative.
		
		Consequently,
		
		\[
		\boxed{
			E[X]\ge0.
		}
		\]
		
		\bigskip
		
		\paragraph{Example 1}
		
		Suppose
		
		\[
		X=
		\begin{cases}
			0,&P=\dfrac14,\\[0.2cm]
			2,&P=\dfrac12,\\[0.2cm]
			6,&P=\dfrac14.
		\end{cases}
		\]
		
		Then
		
		\[
		\begin{aligned}
			E[X]
			&=
			0\left(\frac14\right)
			+
			2\left(\frac12\right)
			+
			6\left(\frac14\right)
			\\[0.3cm]
			&=
			0
			+
			1
			+
			1.5
			\\[0.3cm]
			&=
			2.5.
		\end{aligned}
		\]
		
		As predicted,
		
		\[
		\boxed{
			E[X]=2.5\ge0.
		}
		\]
		
		\bigskip
		
		\paragraph{Example 2}
		
		Suppose
		
		\[
		X\sim U(0,3).
		\]
		
		Its density is
		
		\[
		f_X(x)=
		\frac13,
		\qquad
		0\le x\le3.
		\]
		
		The expectation is
		
		\[
		\begin{aligned}
			E[X]
			&=
			\int_0^3
			x
			\left(\frac13\right)
			dx
			\\[0.3cm]
			&=
			\frac13
			\left[
			\frac{x^2}{2}
			\right]_0^3
			\\[0.3cm]
			&=
			\frac13
			\cdot
			\frac92
			\\[0.3cm]
			&=
			\frac32.
		\end{aligned}
		\]
		
		Again,
		
		\[
		\boxed{
			E[X]=\frac32>0.
		}
		\]
		
		\bigskip
		
		\paragraph{Important Consequences}
		
		The non-negativity property immediately implies several useful results.
		
		If
		
		\[
		X^2\ge0,
		\]
		
		then
		
		\[
		\boxed{
			E[X^2]\ge0.
		}
		\]
		
		Likewise,
		
		\[
		|X|\ge0
		\]
		
		implies
		
		\[
		\boxed{
			E[|X|]\ge0.
		}
		\]
		
		Similarly,
		
		\[
		(X-\mu)^2\ge0,
		\]
		
		where
		
		\[
		\mu=E[X],
		\]
		
		gives
		
		\[
		E[(X-\mu)^2]\ge0.
		\]
		
		Since
		
		\[
		\operatorname{Var}(X)
		=
		E[(X-\mu)^2],
		\]
		
		we immediately conclude
		
		\[
		\boxed{
			\operatorname{Var}(X)\ge0.
		}
		\]
		
		Thus, one of the most important properties of variance follows directly from
		the non-negativity of expectation.
		
		\bigskip
		
		\paragraph{Importance of the Result}
		
		Although simple, the non-negativity property is fundamental throughout
		probability and statistics.
		
		It is the starting point for many important inequalities, including
		
		\[
		\operatorname{Var}(X)\ge0,
		\]
		
		Markov's Inequality,
		
		Chebyshev's Inequality,
		
		Jensen's Inequality,
		
		and numerous other results involving expectations of nonnegative random
		variables.
		
		Together with linearity and monotonicity, it forms one of the basic axioms
		governing the behavior of the expectation operator.
		
	\end{mdframed}
	
	\section{Absolute Value Inequality}
	\begin{mdframed}[
		linewidth=0.8pt,
		roundcorner=6pt,
		backgroundcolor=gray!5,
		linecolor=gray!60,
		innertopmargin=10pt,
		innerbottommargin=10pt,
		innerleftmargin=10pt,
		innerrightmargin=10pt
		]
		
		\subsection*{Absolute Value Inequality}
		
		One of the most useful inequalities involving expected values states that the
		absolute value of the expectation can never exceed the expectation of the
		absolute value.
		
		This inequality plays a central role in probability theory, mathematical
		statistics, and analysis. It is frequently used to bound expectations whose
		exact values are difficult to compute.
		
		\bigskip
		
		\paragraph{Theorem (Absolute Value Inequality)}
		
		Let \(X\) be a random variable such that
		
		\[
		E[|X|]
		<
		\infty.
		\]
		
		Then
		
		\[
		\boxed{
			|E[X]|
			\le
			E[|X|].
		}
		\]
		
		This inequality is sometimes called the \emph{triangle inequality for
			expectation}.
		
		\bigskip
		
		\paragraph{Interpretation}
		
		The quantity
		
		\[
		E[X]
		\]
		
		is the average value of the random variable.
		
		Positive and negative values of \(X\) may cancel one another, causing the
		expectation to be relatively small.
		
		On the other hand,
		
		\[
		E[|X|]
		\]
		
		measures the average magnitude of the random variable.
		
		Since taking absolute values removes all cancellations, it is natural that
		
		\[
		E[|X|]
		\]
		
		is always at least as large as
		
		\[
		|E[X]|.
		\]
		
		\bigskip
		
		\paragraph{Proof}
		
		For every real number \(x\),
		
		\[
		-|x|
		\le
		x
		\le
		|x|.
		\]
		
		Replacing \(x\) by the random variable \(X\),
		
		\[
		-|X|
		\le
		X
		\le
		|X|.
		\]
		
		Applying the monotonicity of expectation,
		
		\[
		E[-|X|]
		\le
		E[X]
		\le
		E[|X|].
		\]
		
		Using the linearity of expectation,
		
		\[
		E[-|X|]
		=
		-\,E[|X|].
		\]
		
		Therefore,
		
		\[
		-\,E[|X|]
		\le
		E[X]
		\le
		E[|X|].
		\]
		
		Thus, the expectation lies inside the interval
		
		\[
		[-E[|X|],\,E[|X|]].
		\]
		
		A basic property of real numbers states that
		
		\[
		-a
		\le
		b
		\le
		a
		\qquad
		(a\ge0)
		\]
		
		if and only if
		
		\[
		|b|
		\le
		a.
		\]
		
		Applying this property with
		
		\[
		a=E[|X|]
		\]
		
		and
		
		\[
		b=E[X],
		\]
		
		we conclude
		
		\[
		\boxed{
			|E[X]|
			\le
			E[|X|].
		}
		\]
		
		This completes the proof.
		
		\bigskip
		
		\paragraph{Alternative Proof (Discrete Case)}
		
		Suppose that \(X\) is a discrete random variable.
		
		Its expectation is
		
		\[
		E[X]
		=
		\sum_x
		xP(X=x).
		\]
		
		Taking absolute values,
		
		\[
		\left|
		E[X]
		\right|
		=
		\left|
		\sum_x
		xP(X=x)
		\right|.
		\]
		
		Applying the triangle inequality for finite or infinite sums,
		
		\[
		\left|
		\sum_x
		xP(X=x)
		\right|
		\le
		\sum_x
		|x|P(X=x).
		\]
		
		Recognizing the last summation,
		
		\[
		\boxed{
			|E[X]|
			\le
			E[|X|].
		}
		\]
		
		\bigskip
		
		\paragraph{Alternative Proof (Continuous Case)}
		
		Suppose that \(X\) has probability density function
		
		\[
		f_X(x).
		\]
		
		Then
		
		\[
		E[X]
		=
		\int_{-\infty}^{\infty}
		x
		f_X(x)
		\,dx.
		\]
		
		Taking absolute values,
		
		\[
		\left|
		E[X]
		\right|
		=
		\left|
		\int_{-\infty}^{\infty}
		x
		f_X(x)
		\,dx
		\right|.
		\]
		
		Using the triangle inequality for integrals,
		
		\[
		\left|
		\int
		g(x)\,dx
		\right|
		\le
		\int
		|g(x)|\,dx,
		\]
		
		with
		
		\[
		g(x)=xf_X(x),
		\]
		
		we obtain
		
		\[
		\begin{aligned}
			|E[X]|
			&\le
			\int_{-\infty}^{\infty}
			|x|
			f_X(x)
			\,dx
			\\[0.3cm]
			&=
			E[|X|].
		\end{aligned}
		\]
		
		Therefore,
		
		\[
		\boxed{
			|E[X]|
			\le
			E[|X|].
		}
		\]
		
		\bigskip
		
		\paragraph{Example 1}
		
		Suppose
		
		\[
		X=
		\begin{cases}
			-2,&P=\dfrac12,\\[0.2cm]
			4,&P=\dfrac12.
		\end{cases}
		\]
		
		First,
		
		\[
		\begin{aligned}
			E[X]
			&=
			(-2)\left(\frac12\right)
			+
			4\left(\frac12\right)
			\\[0.3cm]
			&=
			-1+2
			\\[0.3cm]
			&=
			1.
		\end{aligned}
		\]
		
		Hence,
		
		\[
		|E[X]|=1.
		\]
		
		Now compute
		
		\[
		E[|X|].
		\]
		
		Since
		
		\[
		|X|=
		\begin{cases}
			2,&P=\dfrac12,\\[0.2cm]
			4,&P=\dfrac12,
		\end{cases}
		\]
		
		we have
		
		\[
		\begin{aligned}
			E[|X|]
			&=
			2\left(\frac12\right)
			+
			4\left(\frac12\right)
			\\[0.3cm]
			&=
			1+2
			\\[0.3cm]
			&=
			3.
		\end{aligned}
		\]
		
		Therefore,
		
		\[
		|E[X]|
		=
		1
		<
		3
		=
		E[|X|].
		\]
		
		\bigskip
		
		\paragraph{Example 2}
		
		Suppose
		
		\[
		X=
		\begin{cases}
			-1,&P=\dfrac12,\\[0.2cm]
			1,&P=\dfrac12.
		\end{cases}
		\]
		
		Then
		
		\[
		E[X]
		=
		0,
		\]
		
		while
		
		\[
		E[|X|]
		=
		1.
		\]
		
		Thus,
		
		\[
		\boxed{
			0
			=
			|E[X]|
			<
			E[|X|]
			=
			1.
		}
		\]
		
		This example illustrates how positive and negative values may cancel in the
		expectation but not in the expectation of the absolute value.
		
		\bigskip
		
		\paragraph{When Does Equality Hold?}
		
		Equality holds if and only if the random variable never changes sign.
		
		More precisely,
		
		\[
		|E[X]|
		=
		E[|X|]
		\]
		
		if and only if either
		
		\[
		X\ge0
		\quad
		\text{almost surely},
		\]
		
		or
		
		\[
		X\le0
		\quad
		\text{almost surely}.
		\]
		
		Indeed, if \(X\) is always nonnegative,
		
		\[
		|X|=X,
		\]
		
		so
		
		\[
		|E[X]|
		=
		E[X]
		=
		E[|X|].
		\]
		
		Similarly, if \(X\) is always nonpositive,
		
		\[
		|X|=-X,
		\]
		
		and
		
		\[
		|E[X]|
		=
		-E[X]
		=
		E[|X|].
		\]
		
		Whenever the random variable takes both positive and negative values with
		positive probability, strict inequality generally holds because of
		cancellation.
		
		\bigskip
		
		\paragraph{Importance of the Result}
		
		The Absolute Value Inequality is one of the fundamental inequalities of
		probability theory.
		
		It is frequently used to establish upper bounds on expectations and serves
		as the foundation for many deeper inequalities in analysis and probability,
		including Hölder's Inequality, Minkowski's Inequality, Jensen's Inequality,
		and various convergence theorems.
		
		Together with linearity, monotonicity, and non-negativity, it forms one of
		the essential algebraic properties of the expectation operator.
		
	\end{mdframed}
	
	\section{Indicator Random Variables}
	\begin{mdframed}[
		linewidth=0.8pt,
		roundcorner=6pt,
		backgroundcolor=gray!5,
		linecolor=gray!60,
		innertopmargin=10pt,
		innerbottommargin=10pt,
		innerleftmargin=10pt,
		innerrightmargin=10pt
		]
		
		\subsection*{Indicator Random Variables}
		
		Indicator random variables are among the most useful tools in probability and
		statistics. They transform logical statements or events into random
		variables that take only the values \(0\) and \(1\).
		
		Their simplicity makes them extremely powerful, especially when computing
		expectations, probabilities, and counting random quantities.
		
		Many elegant probabilistic proofs rely almost entirely on indicator random
		variables and the linearity of expectation.
		
		\bigskip
		
		\paragraph{Definition}
		
		Let
		
		\[
		A
		\]
		
		be an event in a probability space.
		
		The \emph{indicator random variable} associated with the event \(A\) is
		defined by
		
		\[
		\boxed{
			I_A=
			\begin{cases}
				1,&\text{if }A\text{ occurs},\\[0.2cm]
				0,&\text{if }A\text{ does not occur}.
			\end{cases}
		}
		\]
		
		Thus, an indicator variable simply records whether an event has occurred.
		
		If the event happens,
		
		\[
		I_A=1,
		\]
		
		otherwise,
		
		\[
		I_A=0.
		\]
		
		\bigskip
		
		\paragraph{Probability Distribution}
		
		Since the indicator variable takes only two possible values,
		
		\[
		0
		\quad\text{and}\quad
		1,
		\]
		
		its probability distribution is particularly simple.
		
		Indeed,
		
		\[
		P(I_A=1)=P(A),
		\]
		
		and
		
		\[
		P(I_A=0)=1-P(A).
		\]
		
		Therefore,
		
		\[
		I_A
		\sim
		\text{Bernoulli}(P(A)).
		\]
		
		\bigskip
		
		\paragraph{Expectation of an Indicator Variable}
		
		The expectation follows directly from the definition.
		
		Using the discrete expectation formula,
		
		\[
		E[I_A]
		=
		\sum_x
		xP(I_A=x).
		\]
		
		Since the indicator takes only the values \(0\) and \(1\),
		
		\[
		\begin{aligned}
			E[I_A]
			&=
			0\cdot P(I_A=0)
			+
			1\cdot P(I_A=1)
			\\[0.3cm]
			&=
			0\cdot(1-P(A))
			+
			1\cdot P(A)
			\\[0.3cm]
			&=
			P(A).
		\end{aligned}
		\]
		
		Hence,
		
		\[
		\boxed{
			E[I_A]=P(A).
		}
		\]
		
		This is one of the most important identities in probability theory.
		
		It allows probabilities to be computed using expectations.
		
		\bigskip
		
		\paragraph{Variance of an Indicator Variable}
		
		Since
		
		\[
		I_A^2=I_A,
		\]
		
		we have
		
		\[
		E[I_A^2]
		=
		E[I_A]
		=
		P(A).
		\]
		
		Using the variance formula,
		
		\[
		\operatorname{Var}(X)
		=
		E[X^2]-(E[X])^2,
		\]
		
		we obtain
		
		\[
		\begin{aligned}
			\operatorname{Var}(I_A)
			&=
			P(A)-P(A)^2
			\\[0.3cm]
			&=
			P(A)(1-P(A)).
		\end{aligned}
		\]
		
		Thus,
		
		\[
		\boxed{
			\operatorname{Var}(I_A)
			=
			P(A)(1-P(A)).
		}
		\]
		
		This is precisely the variance of a Bernoulli random variable.
		
		\bigskip
		
		\paragraph{Indicator Variables and Counting}
		
		Indicator variables become especially useful when counting objects.
		
		Suppose that
		
		\[
		A_1,A_2,\ldots,A_n
		\]
		
		are events.
		
		Define
		
		\[
		I_i=
		I_{A_i},
		\qquad
		i=1,\ldots,n.
		\]
		
		Now consider
		
		\[
		N
		=
		I_1+I_2+\cdots+I_n.
		\]
		
		Since each indicator contributes either \(0\) or \(1\),
		
		\[
		N
		\]
		
		counts the number of events that occur.
		
		\bigskip
		
		\paragraph{Expected Number of Occurring Events}
		
		Applying the linearity of expectation,
		
		\[
		\begin{aligned}
			E[N]
			&=
			E[I_1+\cdots+I_n]
			\\[0.3cm]
			&=
			E[I_1]+\cdots+E[I_n].
		\end{aligned}
		\]
		
		Using
		
		\[
		E[I_i]=P(A_i),
		\]
		
		we obtain
		
		\[
		\boxed{
			E[N]
			=
			\sum_{i=1}^{n}
			P(A_i).
		}
		\]
		
		Notice that no independence assumption is required.
		
		This is one of the most useful consequences of the linearity of expectation.
		
		\bigskip
		
		\paragraph{Example 1: Tossing a Coin}
		
		Suppose a fair coin is tossed once.
		
		Define
		
		\[
		I=
		\begin{cases}
			1,&\text{if Heads occurs},\\
			0,&\text{otherwise}.
		\end{cases}
		\]
		
		Since
		
		\[
		P(\text{Heads})=\frac12,
		\]
		
		we have
		
		\[
		\boxed{
			E[I]
			=
			\frac12.
		}
		\]
		
		The expected value of the indicator is exactly the probability of obtaining
		Heads.
		
		\bigskip
		
		\paragraph{Example 2: Number of Heads in Five Tosses}
		
		Suppose a fair coin is tossed five times.
		
		For each toss, define
		
		\[
		I_i=
		\begin{cases}
			1,&\text{if toss }i\text{ is Heads},\\
			0,&\text{otherwise}.
		\end{cases}
		\]
		
		The total number of Heads is
		
		\[
		H
		=
		I_1+I_2+I_3+I_4+I_5.
		\]
		
		Applying linearity,
		
		\[
		\begin{aligned}
			E[H]
			&=
			E[I_1]+\cdots+E[I_5]
			\\[0.3cm]
			&=
			5\left(\frac12\right)
			\\[0.3cm]
			&=
			2.5.
		\end{aligned}
		\]
		
		Thus,
		
		\[
		\boxed{
			E[H]=2.5.
		}
		\]
		
		Notice that we never needed the distribution of \(H\), which is actually
		binomial.
		
		\bigskip
		
		\paragraph{Example 3: Expected Number of Fixed Points}
		
		Suppose a permutation of
		
		\[
		\{1,2,\ldots,n\}
		\]
		
		is chosen uniformly at random.
		
		Let
		
		\[
		I_i=
		\begin{cases}
			1,&\text{if }i\text{ remains in its original position},\\
			0,&\text{otherwise}.
		\end{cases}
		\]
		
		The total number of fixed points is
		
		\[
		F
		=
		\sum_{i=1}^{n}I_i.
		\]
		
		Since
		
		\[
		P(i\text{ is fixed})
		=
		\frac1n,
		\]
		
		we have
		
		\[
		E[I_i]
		=
		\frac1n.
		\]
		
		Therefore,
		
		\[
		\begin{aligned}
			E[F]
			&=
			\sum_{i=1}^{n}
			E[I_i]
			\\[0.3cm]
			&=
			n\left(\frac1n\right)
			\\[0.3cm]
			&=
			1.
		\end{aligned}
		\]
		
		Hence,
		
		\[
		\boxed{
			E[F]=1.
		}
		\]
		
		Remarkably, this result holds for every positive integer \(n\).
		
		\bigskip
		
		\paragraph{Importance of Indicator Variables}
		
		Indicator random variables provide a bridge between probabilities and
		expectations through the identity
		
		\[
		\boxed{
			E[I_A]=P(A).
		}
		\]
		
		Combined with the linearity of expectation, this simple identity allows many
		counting problems to be solved without computing complicated probability
		distributions.
		
		Indicator variables are fundamental in combinatorics, graph theory,
		stochastic processes, randomized algorithms, statistical inference, and many
		other areas of probability.
		
		For this reason, they are considered one of the most powerful tools in
		elementary probability theory.
		
	\end{mdframed}
	
	\section{Sample Mean is Unbiased}
	
	\begin{mdframed}[
		linewidth=0.8pt,
		roundcorner=6pt,
		backgroundcolor=gray!5,
		linecolor=gray!60,
		innertopmargin=10pt,
		innerbottommargin=10pt,
		innerleftmargin=10pt,
		innerrightmargin=10pt
		]
		
		\subsection*{The Sample Mean is an Unbiased Estimator}
		
		One of the most fundamental results in statistics is that the sample mean is
		an \emph{unbiased estimator} of the population mean.
		
		This property explains why the sample mean is the natural estimator of the
		unknown population mean. On average, it neither systematically overestimates
		nor underestimates the true parameter.
		
		\bigskip
		
		\paragraph{Definition of the Sample Mean}
		
		Suppose
		
		\[
		X_1,X_2,\ldots,X_n
		\]
		
		are independent and identically distributed (i.i.d.) random variables with
		
		\[
		E[X_i]=\mu,
		\qquad
		i=1,\ldots,n.
		\]
		
		The sample mean is defined as
		
		\[
		\boxed{
			\bar X
			=
			\frac1n
			\sum_{i=1}^{n}X_i.
		}
		\]
		
		Its purpose is to estimate the unknown population mean
		
		\[
		\mu.
		\]
		
		\bigskip
		
		\paragraph{Definition of an Unbiased Estimator}
		
		An estimator
		
		\[
		\hat\theta
		\]
		
		of a parameter
		
		\[
		\theta
		\]
		
		is said to be \emph{unbiased} if
		
		\[
		\boxed{
			E[\hat\theta]
			=
			\theta.
		}
		\]
		
		In other words, the average value of the estimator over repeated sampling is
		exactly the parameter being estimated.
		
		Thus, to prove that the sample mean is unbiased, we must show that
		
		\[
		E[\bar X]
		=
		\mu.
		\]
		
		\bigskip
		
		\paragraph{Step 1: Start from the Definition}
		
		By definition,
		
		\[
		\bar X
		=
		\frac1n
		\sum_{i=1}^{n}X_i.
		\]
		
		Taking expectations on both sides,
		
		\[
		E[\bar X]
		=
		E\!\left(
		\frac1n
		\sum_{i=1}^{n}X_i
		\right).
		\]
		
		\bigskip
		
		\paragraph{Step 2: Apply the Constant Multiple Rule}
		
		Since
		
		\[
		\frac1n
		\]
		
		is a constant, it can be taken outside the expectation,
		
		\[
		\begin{aligned}
			E[\bar X]
			&=
			\frac1n
			E\!\left(
			\sum_{i=1}^{n}X_i
			\right).
		\end{aligned}
		\]
		
		\bigskip
		
		\paragraph{Step 3: Apply the Linearity of Expectation}
		
		Using the linearity of expectation,
		
		\[
		E\!\left(
		\sum_{i=1}^{n}X_i
		\right)
		=
		\sum_{i=1}^{n}E[X_i].
		\]
		
		Therefore,
		
		\[
		\begin{aligned}
			E[\bar X]
			&=
			\frac1n
			\sum_{i=1}^{n}E[X_i].
		\end{aligned}
		\]
		
		Notice that this step does \textbf{not} require the random variables to be
		independent. Linearity of expectation always holds.
		
		\bigskip
		
		\paragraph{Step 4: Use Identical Distribution}
		
		Since the observations are identically distributed,
		
		\[
		E[X_i]=\mu
		\qquad
		(i=1,\ldots,n).
		\]
		
		Substituting,
		
		\[
		\begin{aligned}
			E[\bar X]
			&=
			\frac1n
			(\mu+\mu+\cdots+\mu).
		\end{aligned}
		\]
		
		There are exactly \(n\) copies of \(\mu\), so
		
		\[
		\mu+\mu+\cdots+\mu
		=
		n\mu.
		\]
		
		Hence,
		
		\[
		\begin{aligned}
			E[\bar X]
			&=
			\frac1n(n\mu)
			\\[0.3cm]
			&=
			\mu.
		\end{aligned}
		\]
		
		Therefore,
		
		\[
		\boxed{
			E[\bar X]=\mu.
		}
		\]
		
		This proves that the sample mean is an unbiased estimator of the population
		mean.
		
		\bigskip
		
		\paragraph{Interpretation}
		
		Imagine repeatedly drawing random samples of size \(n\) from the same
		population.
		
		Each sample produces a different sample mean,
		
		\[
		\bar X_1,\,
		\bar X_2,\,
		\bar X_3,\,
		\ldots
		\]
		
		Some sample means will be larger than the true population mean, while others
		will be smaller.
		
		However, if we average all these sample means, we obtain exactly
		
		\[
		\mu.
		\]
		
		Thus, the estimator has no systematic tendency to overestimate or
		underestimate the parameter.
		
		\bigskip
		
		\paragraph{Example}
		
		Suppose a population has mean
		
		\[
		\mu=50.
		\]
		
		We repeatedly draw random samples of size
		
		\[
		n=10.
		\]
		
		Although the sample means may vary,
		
		\[
		48.7,\quad
		52.1,\quad
		49.5,\quad
		50.9,\quad
		47.8,\quad
		\ldots,
		\]
		
		their long-run average is
		
		\[
		50.
		\]
		
		Therefore,
		
		\[
		E[\bar X]=50.
		\]
		
		The estimator fluctuates because of sampling variability, but it is centered
		at the true population mean.
		
		\bigskip
		
		\paragraph{Why Independence Is Mentioned}
		
		The proof itself does not require independence.
		
		Only the assumptions
		
		\[
		E[X_i]=\mu
		\]
		
		for all \(i\) are needed.
		
		Independence becomes important when studying the variance of the sample
		mean,
		
		\[
		\operatorname{Var}(\bar X),
		\]
		
		or when deriving its sampling distribution, such as in the Central Limit
		Theorem.
		
		Thus,
		
		\[
		\boxed{
			\text{Unbiasedness requires only identical means,}
		}
		\]
		
		whereas
		
		\[
		\boxed{
			\text{Variance formulas and sampling distributions require additional assumptions,}
		}
		\]
		
		typically independence.
		
		\bigskip
		
		\paragraph{Importance of the Result}
		
		The unbiasedness of the sample mean is one of the cornerstones of
		mathematical statistics.
		
		It justifies using
		
		\[
		\bar X
		\]
		
		as an estimator of the unknown population mean and serves as the foundation
		for confidence intervals, hypothesis testing, regression analysis, the
		Central Limit Theorem, and many other statistical procedures.
		
		Together with its variance,
		
		\[
		\operatorname{Var}(\bar X)=\frac{\sigma^2}{n},
		\]
		
		this result explains why the sample mean becomes increasingly accurate as
		the sample size grows.
		
	\end{mdframed}
	
	
	\section{Summary Table}
	\begin{mdframed}[
		linewidth=0.8pt,
		roundcorner=6pt,
		backgroundcolor=gray!5,
		linecolor=gray!60,
		innertopmargin=10pt,
		innerbottommargin=10pt,
		innerleftmargin=10pt,
		innerrightmargin=10pt
		]
		
		\subsection*{Summary Table}
		
		The table below summarizes the most important properties of the expectation
		operator presented in this chapter.
		
		\bigskip
		
		\renewcommand{\arraystretch}{1.5}
		
		\begin{center}
			\begin{tabular}{|p{5.2cm}|p{7.8cm}|}
				\hline
				\textbf{Property} &
				\textbf{Formula}
				\\
				\hline
				
				Definition (Discrete)
				&
				\[
				E[X]
				=
				\sum_x xP(X=x)
				\]
				\\
				\hline
				
				Definition (Continuous)
				&
				\[
				E[X]
				=
				\int_{-\infty}^{\infty}
				xf_X(x)\,dx
				\]
				\\
				\hline
				
				Expectation of a Function (LOTUS)
				&
				\[
				E[g(X)]
				=
				\sum_x g(x)P(X=x)
				\]
				
				or
				
				\[
				E[g(X)]
				=
				\int_{-\infty}^{\infty}
				g(x)f_X(x)\,dx
				\]
				\\
				\hline
				
				Expectation of a Constant
				&
				\[
				E[c]=c
				\]
				\\
				\hline
				
				Linearity of Expectation
				&
				\[
				E[aX+bY]
				=
				aE[X]
				+
				bE[Y]
				\]
				\\
				\hline
				
				Expectation of Sums
				&
				\[
				E\!\left(
				\sum_{i=1}^{n}X_i
				\right)
				=
				\sum_{i=1}^{n}E[X_i]
				\]
				\\
				\hline
				
				Expectation of Products
				(Independent Variables)
				&
				\[
				E[XY]
				=
				E[X]E[Y]
				\]
				\\
				\hline
				
				Monotonicity
				&
				\[
				X\le Y
				\Longrightarrow
				E[X]\le E[Y]
				\]
				\\
				\hline
				
				Non-negativity
				&
				\[
				X\ge0
				\Longrightarrow
				E[X]\ge0
				\]
				\\
				\hline
				
				Absolute Value Inequality
				&
				\[
				|E[X]|
				\le
				E[|X|]
				\]
				\\
				\hline
				
				Indicator Variables
				&
				\[
				E[I_A]
				=
				P(A)
				\]
				\\
				\hline
				
				Sample Mean is Unbiased
				&
				\[
				E[\bar X]
				=
				\mu
				\]
				\\
				\hline
				
			\end{tabular}
		\end{center}
		
		\bigskip
		
		These properties constitute the algebraic foundation of the expectation
		operator and are among the most frequently used results in probability,
		statistics, stochastic processes, machine learning, econometrics, and
		statistical inference.
		
		Most advanced results involving expectations—including variance,
		covariance, moment generating functions, laws of large numbers, the Central
		Limit Theorem, and estimation theory—are direct consequences or extensions
		of these fundamental properties.
		
	\end{mdframed}
	
	
\end{document}